\chapter{Parathyroid Hormone (PTH)}

\section*{What Is This Test?}
Parathyroid hormone (PTH, also called parathormone) is an 84-amino-acid single-chain polypeptide hormone produced and secreted by the chief cells of the four parathyroid glands located on the posterior aspect of the thyroid gland. PTH is the primary regulator of serum calcium and phosphate homeostasis. PTH secretion is regulated by the calcium-sensing receptor (CaSR) on parathyroid chief cells: a decrease in ionized calcium stimulates PTH secretion; an increase in ionized calcium suppresses PTH secretion. PTH acts on three target organs: (1) bone --- stimulates osteoclast-mediated bone resorption releasing calcium and phosphate; (2) kidneys --- increases distal tubular calcium reabsorption, decreases proximal tubular phosphate reabsorption (phosphaturia), and stimulates 1-alpha-hydroxylase (CYP27B1) to convert 25-hydroxyvitamin D to active 1,25-dihydroxyvitamin D (calcitriol); and (3) intestine --- indirectly, via calcitriol, increases intestinal calcium and phosphate absorption. The net effect is to raise serum calcium and lower serum phosphate. PTH has a very short plasma half-life of 2--4 minutes and is cleaved in the liver and kidneys into N-terminal, C-terminal, and mid-region fragments. The intact PTH assay (iPTH) measures the biologically active full-length 1-84 molecule and is the standard assay for clinical use. PTH must always be interpreted with a simultaneously drawn serum calcium, phosphate, albumin, and 25-hydroxyvitamin D level \cite{tietz2012textbook}.

\section*{Alternative Names}
Intact PTH (iPTH); Parathormone; Whole PTH (1-84); Bioactive PTH; PTH, Intact Molecule; Serum PTH; Plasma PTH; PTH-Related Peptide (PTHrP, distinct analyte)

\section*{Principle}
Intact PTH (iPTH) is measured using a two-site sandwich immunometric (chemiluminescent) immunoassay. Two antibodies are used: a capture antibody directed against the C-terminal or mid-region (amino acids 39--84) and a detection antibody directed against the N-terminal region (amino acids 1--34). This design ensures that only the full-length biologically active PTH (1-84) is measured and N-terminal truncated fragments (such as PTH 7-84, which accumulates in renal failure and antagonizes the PTH receptor) are not detected. Common platforms: Roche Cobas ECLIA, Abbott Architect CMIA, Siemens Atellica/ADVIA Centaur CLIA, and Beckman Coulter DxI. Third-generation assays use antibodies against the very N-terminus (amino acids 1--4) to exclude all N-terminally truncated fragments, and are sometimes called "whole PTH" or "bioactive PTH" assays. However, second-generation intact PTH assays remain the most widely used. The functional sensitivity is typically <3--5 pg/mL. PTH is unstable at room temperature; specimens should be centrifuged within 2 hours and serum separated and refrigerated or frozen. Pre-analytical stability is critical: PTH half-life in vitro at room temperature is approximately 4--6 hours.

\section*{History}
The parathyroid glands were first described anatomically by Richard Owen in 1850 and their function was discovered by Ivar Sandstrom in 1880. The first parathyroidectomy was performed by Felix Mandl in Vienna in 1925 on a patient with osteitis fibrosa cystica caused by a parathyroid adenoma. PTH was first isolated by Collip in 1925. The first radioimmunoassay for PTH was developed by Berson and Yalow in 1963 (the same team that developed the insulin RIA). The first intact PTH immunometric assay was developed by Nussbaum and colleagues in 1987 \cite{nussbaum1987highly}. The development of the calcium-sensing receptor (CaSR) concept and its cloning by Brown and colleagues in 1993 revolutionized understanding of PTH regulation and led to the development of calcimimetics (cinacalcet) for secondary hyperparathyroidism \cite{brown1993cloning}. The 2014 guidelines for primary hyperparathyroidism (Bilezikian et al.) and 2016 guidelines for hypoparathyroidism (Brandi et al.) standardized PTH interpretation and clinical management \cite{bilezikian2014guidelines,brandi2016management}. The rapid intraoperative PTH assay (allowing 15--20 minute turnaround) revolutionized parathyroid surgery by enabling confirmation of adenoma removal during the operation.

\section*{Why This Test Is Ordered}
\begin{itemize}
  \item Differential diagnosis of hypercalcemia: elevated or inappropriately normal PTH with hypercalcemia indicates primary hyperparathyroidism; suppressed PTH (<20 pg/mL) with hypercalcemia indicates non-PTH-mediated hypercalcemia (malignancy, granulomatous disease, milk-alkali syndrome)
  \item Diagnosis of primary hyperparathyroidism: elevated PTH (>65 pg/mL) with elevated or high-normal serum calcium (>10.2 mg/dL) and low serum phosphate (<2.5 mg/dL). The most common cause is a solitary parathyroid adenoma (80--85\%)
  \item Diagnosis of hypoparathyroidism: low or inappropriately normal PTH with hypocalcemia (<8.5 mg/dL) and hyperphosphatemia, most commonly after thyroid/parathyroid surgery
  \item Management of secondary hyperparathyroidism in chronic kidney disease (CKD) --- target PTH 2--9 times the upper limit of normal for dialysis patients per KDIGO guidelines \cite{kdigo2017clinical}
  \item Intraoperative PTH monitoring during parathyroidectomy for primary hyperparathyroidism --- a drop in PTH >50\% from pre-excision baseline at 10 minutes after adenoma removal confirms biochemical cure
  \item Evaluation of familial hypocalciuric hypercalcemia (FHH) vs. primary hyperparathyroidism --- low 24-hour urinary calcium clearance/creatinine clearance ratio (<0.01) suggests FHH
  \item Investigation of osteoporosis --- hyperparathyroidism causes preferential cortical bone loss
  \item Diagnosis of tertiary hyperparathyroidism in renal transplant recipients with persistent hypercalcemia
  \item Assessment of vitamin D deficiency --- secondary hyperparathyroidism with normal or low-normal serum calcium and low 25-hydroxyvitamin D
\end{itemize}

\section*{Primary vs Secondary Test}
PTH is a primary test in the differential diagnosis of hypercalcemia and hypocalcemia. It is the single most informative test for classifying disorders of calcium metabolism. PTH must always be interpreted with a simultaneously drawn serum calcium (total, corrected for albumin, or ionized), phosphate, and 25-hydroxyvitamin D.

\section*{How This Test Is Performed}
A healthcare professional draws a blood sample from a vein into a serum separator tube (gold-top) or plain red-top tube (~3--5 mL). The specimen should be centrifuged within 2 hours and serum separated promptly, because PTH degrades at room temperature (in vitro half-life ~4--6 hours). If processing is delayed, the specimen should be refrigerated. For intra-operative PTH monitoring, blood is drawn before incision (baseline) and at 5, 10, and 15 minutes after removal of the suspected adenoma. The sample should be drawn from a site without IV calcium or phosphate infusion.

\section*{What Preparation Is Needed}
Fasting is recommended (8--10 hours) because food intake can modestly alter serum calcium and PTH. A simultaneously drawn serum calcium (total plus albumin, or ionized) and 25-hydroxyvitamin D are essential. High-dose biotin (>5 mg/day) must be held 48 hours before testing. The patient should not be receiving IV calcium, phosphate, or calcitonin at the time of blood draw, as these acutely alter PTH. The patient's posture should be considered: PTH is slightly higher in the supine than upright position (possibly due to hemodynamic effects). For optimal consistency, the blood draw should be done at approximately the same time of day for serial measurements, as PTH has a modest circadian rhythm (higher at night, lower in late morning, approximately 20--30\% variation).

\section*{Contraindications and Precautions}
\begin{itemize}
  \item PTH MUST be interpreted with a simultaneously drawn serum calcium, phosphate, and 25-hydroxyvitamin D. Isolated PTH elevation without these values is uninterpretable. PTH degrades in vitro at room temperature: specimens left at room temperature >8 hours may yield falsely low results. Hemolysis interferes with some assays. Severe lipemia causes optical interference. Biotin interferes with streptavidin-biotin immunoassays. PTH fragments (7-84) accumulate in renal failure (CKD stages 4--5) because of decreased renal clearance. Second-generation intact PTH assays cross-react with PTH 7-84 (approximately 30--50\% cross-reactivity), leading to overestimation of biologically active PTH in CKD. Third-generation "whole PTH" or "bioactive PTH" assays are more specific for the 1-84 molecule and may be preferred in CKD. PTH is highest at night and lowest in the late morning, with a difference of approximately 20--30\%. Drugs that alter PTH: lithium (elevates the calcium set-point for PTH suppression, causing mild hypercalcemia and elevated PTH in 10--20\% of users), thiazide diuretics (decrease urinary calcium excretion, mild hypercalcemia, and lower PTH), loop diuretics (increase calciuria, stimulate PTH), bisphosphonates and denosumab (lower calcium, stimulate PTH), cinacalcet (calcimimetic, suppresses PTH), calcitriol and vitamin D analogs (suppress PTH), and phosphate binders (indirectly affect PTH via phosphate and calcium changes). PTH-related peptide (PTHrP) does not cross-react in PTH assays but causes humoral hypercalcemia of malignancy by activating the PTH receptor.
\end{itemize}
\section*{Risks}
Minimal risk. Slight pain, bruising, or hematoma at venipuncture site. Rarely: vasovagal syncope. For intra-operative PTH monitoring: additional blood draws during surgery, no additional risk beyond the surgical procedure itself.

\section*{What Happens After the Test}
PTH results available within 1--4 hours (routine immunoassay) or 15--20 minutes (intra-operative assay). Interpretation of hypercalcemia: elevated or high-normal PTH (>20 pg/mL) with hypercalcemia indicates PTH-mediated hypercalcemia, most commonly primary hyperparathyroidism. Suppressed PTH (<10--20 pg/mL) with hypercalcemia indicates non-PTH-mediated hypercalcemia, warranting workup for malignancy (PTHrP, bone metastases, myeloma), granulomatous disease (sarcoidosis, tuberculosis), vitamin D intoxication, milk-alkali syndrome, or thyrotoxicosis. In primary hyperparathyroidism, the next step is 24-hour urinary calcium to exclude FHH, and localization with parathyroid ultrasound and sestamibi (Tc-99m MIBI) scan. Surgical referral for parathyroidectomy is indicated if the patient meets criteria: age <50, serum calcium >1 mg/dL above upper limit, osteoporosis (T-score <\-2.5 at any site or vertebral fracture), creatinine clearance <60 mL/min, or nephrolithiasis/nephrocalcinosis.

\section*{Understanding Results}

\subsection*{Below Normal}
\begin{itemize}
  \item Hypoparathyroidism: low or inappropriately normal PTH (typically <10--15 pg/mL) with hypocalcemia (<8.5 mg/dL) and hyperphosphatemia. The most common cause is post-surgical (thyroidectomy, parathyroidectomy) with inadvertent removal or devascularization of parathyroid glands, accounting for ~75\% of cases. Other causes: autoimmune hypoparathyroidism (APS-1, activating CaSR antibodies), magnesium deficiency (severe hypomagnesemia <1.2 mg/dL impairs PTH secretion and causes end-organ resistance), infiltrative diseases (hemochromatosis, Wilson disease, metastases), DiGeorge syndrome (22q11.2 deletion, thymic and parathyroid aplasia), and activating CaSR mutations (autosomal dominant hypocalcemia).
  \item Suppressed PTH with hypercalcemia: PTH-independent hypercalcemia. The intact feedback loop is functioning correctly: hypercalcemia suppresses PTH. Most commonly caused by humoral hypercalcemia of malignancy (PTHrP secretion by solid tumors, especially squamous cell carcinoma of the lung, head/neck, esophagus, and renal cell carcinoma), osteolytic bone metastases (breast cancer, multiple myeloma), granulomatous diseases (sarcoidosis, tuberculosis --- activated macrophages express 1-alpha-hydroxylase, producing excess calcitriol), vitamin D intoxication, milk-alkali syndrome, immobilization, thyrotoxicosis, and thiazide diuretics.
  \item Post-parathyroidectomy: PTH should fall >50\% from baseline within 10 minutes of adenoma removal (intra-operative PTH criteria). Undetectable or low PTH post-operatively indicates surgical cure but can cause transient hypocalcemia (hungry bone syndrome --- rapid bone remineralization consuming calcium; resolves over weeks with calcium and calcitriol supplementation).
\end{itemize}

\subsection*{Above Normal}
\begin{itemize}
  \item Primary hyperparathyroidism: elevated or inappropriately normal PTH (>65 pg/mL, often 100--300 pg/mL) with hypercalcemia (>10.2 mg/dL) and hypophosphatemia. The most common cause is a solitary parathyroid adenoma (80--85\%). Other causes: parathyroid hyperplasia (10--15\%, involving all 4 glands), double adenomas (2--4\%), and parathyroid carcinoma (<1\%, PTH often >3--5 times upper limit of normal, calcium >14 mg/dL). Symptoms: bones (osteoporosis, osteitis fibrosa cystica --- salt-and-pepper skull, brown tumors), stones (nephrolithiasis in 20\%), groans (abdominal pain from constipation, pancreatitis, peptic ulcer disease), psychic overtones (depression, anxiety, cognitive impairment, fatigue). Approximately 80\% of patients with primary hyperparathyroidism are asymptomatic at diagnosis because of routine calcium screening \cite{walker2018primary}.
  \item Secondary hyperparathyroidism: elevated PTH with normal or low serum calcium. Causes: vitamin D deficiency (most common cause worldwide), chronic kidney disease (decreased calcitriol production and phosphate retention), inadequate calcium intake, malabsorption (celiac disease, gastric bypass surgery, chronic pancreatitis), liver disease (decreased 25-hydroxylation of vitamin D), bisphosphonate/denosumab therapy (lower serum calcium, stimulate PTH), and loop diuretics (increase urinary calcium loss). In vitamin D deficiency: PTH 70--150 pg/mL, calcium low-normal, phosphate low, 25(OH)D <20 ng/mL. In CKD: PTH progressively rises as GFR declines below 60 mL/min; target PTH in dialysis is 2--9 times ULN per KDIGO.
  \item Tertiary hyperparathyroidism: autonomous PTH hypersecretion after prolonged secondary hyperparathyroidism, most commonly in renal transplant recipients. Hypercalcemia develops because the hyperplastic parathyroid glands function autonomously and fail to suppress with normal calcium. Often requires parathyroidectomy.
  \item Familial hypocalciuric hypercalcemia (FHH): elevated or normal PTH with mild hypercalcemia, hypocalciuria, and normal renal function. Caused by inactivating CaSR mutations (autosomal dominant). The CaSR set-point is shifted, so higher calcium is required to suppress PTH. Benign condition, does NOT require parathyroidectomy. Distinguished from primary hyperparathyroidism by 24-hour urinary calcium/creatinine clearance ratio <0.01 (low urinary calcium despite hypercalcemia).
  \item Lithium-induced hyperparathyroidism: lithium shifts the CaSR set-point, requiring higher calcium to suppress PTH. Mild hypercalcemia, elevated PTH. Develops in 10--20\% of long-term lithium users. May have parathyroid hyperplasia. Often reversible with lithium discontinuation.
  \item Ectopic PTH secretion: extremely rare. Reported in a handful of cases with lung and ovarian carcinomas producing authentic PTH.
\end{itemize}

\section*{Normal Results}
\begin{itemize}
  \item Intact PTH (adults, eucalcemic, vitamin D sufficient): 10--65 pg/mL (1.1--6.8 pmol/L) --- varies by assay; typical ranges 15--65 pg/mL (Roche, Abbott), 12--72 pg/mL (Siemens)
  \item PTH with normal serum calcium (8.5--10.2 mg/dL): 15--65 pg/mL
  \item PTH in primary hyperparathyroidism: >65 pg/mL (often 100--300 pg/mL) with hypercalcemia; PTH may be inappropriately "normal" (40--65 pg/mL) in the setting of hypercalcemia, which is still abnormal
  \item PTH in hypoparathyroidism: <10--15 pg/mL with hypocalcemia
  \item PTH in vitamin D deficiency: 70--150 pg/mL with normal or low-normal calcium and low 25(OH)D <20 ng/mL
  \item PTH target in CKD Stage 5D (dialysis): 130--585 pg/mL (2--9 times ULN per KDIGO 2017)
  \item Intra-operative PTH: >50\% drop at 10 minutes post-excision from baseline confirms adenoma removal (Miami criteria); some centers use more stringent criteria (drop >50\% AND into normal range)
  \item FHH screen: 24-hour urinary calcium/creatinine clearance ratio <0.01 (suggests FHH); >0.02 (primary hyperparathyroidism)
  \item Note: Reference ranges are assay- and laboratory-specific. Always use laboratory-specific ranges. PTH must be interpreted with total calcium (corrected for albumin), ionized calcium, phosphate, and 25(OH)D.
\end{itemize}

\section*{Follow-up Tests}
\begin{itemize}
  \item Serum calcium (total, with albumin for correction): Must be drawn simultaneously with PTH. Corrected calcium = measured total Ca + 0.8 x (4.0 - albumin). Ionized calcium is preferred when albumin is abnormal or acid-base disorders are present.
  \item Ionized calcium: The biologically active fraction. Not affected by albumin. Preferred in critical illness, CKD, and when total calcium is borderline.
  \item Serum phosphate: PTH lowers phosphate (phosphaturia). Low phosphate supports primary hyperparathyroidism. High phosphate supports hypoparathyroidism.
  \item 25-hydroxyvitamin D: Essential co-test. Vitamin D deficiency is the most common cause of secondary hyperparathyroidism. Target >30 ng/mL.
  \item 1,25-dihydroxyvitamin D (calcitriol): Elevated in primary hyperparathyroidism (PTH-driven conversion). Also elevated in granulomatous disease and lymphoma (extra-renal 1-alpha-hydroxylase).
  \item 24-hour urinary calcium and creatinine: Distinguishes FHH (low urinary calcium, ratio <0.01) from primary hyperparathyroidism (normal or elevated urinary calcium). Also assesses renal stone risk.
  \item Parathyroid ultrasound (high-resolution, 7--15 MHz): First-line localization study. Parathyroid adenoma appears as a hypoechoic, oval mass posterior to the thyroid. Sensitivity 70--85\%.
  \item Tc-99m sestamibi (MIBI) scan with SPECT/CT: Nuclear medicine parathyroid imaging. MIBI is taken up by mitochondria-rich oxyphil cells in parathyroid adenomas and washes out more slowly than from thyroid. Combined with SPECT/CT, sensitivity 85--95\%. Used for pre-operative localization.
  \item 4D-CT parathyroid: Dynamic contrast-enhanced CT with arterial, venous, and delayed phases. Superior sensitivity (90--95\%) when ultrasound and MIBI are negative or discordant.
  \item PTHrP (PTH-related peptide): Measured when hypercalcemia of malignancy is suspected and PTH is suppressed. Elevated PTHrP confirms humoral hypercalcemia of malignancy.
  \item Serum magnesium: Low magnesium impairs PTH secretion and causes end-organ resistance. Must be repleted before diagnosing hypoparathyroidism.
  \item Bone densitometry (DEXA): Assess bone density (cortical bone preferentially affected in hyperparathyroidism --- distal radius, hip).
  \item Renal ultrasound or CT: Assess for nephrolithiasis/nephrocalcinosis in primary hyperparathyroidism.
\end{itemize}

\section*{References}
\bibliographystyle{unsrtnat}
\bibliography{references}

\clearpage
\section*{Regulatory Status and Authorization Date}
This chapter describes a test category, not a specific commercial product. FDA, CDSCO, EU IVDR, and MHRA approval or registration dates therefore cannot be assigned without the assay manufacturer, product name, instrument, and intended use. For laboratory-developed tests, an individual agency approval date may not exist. See the Preface for the interpretation of regulatory status.

\clearpage
